\chapter{Schwerpunkt-Trassierung}

\section{Motivation}

\begin{remark}
Classical railway alignment design separates:
\begin{itemize}
\item geometric design (alignment),
\item cant design,
\item dynamic evaluation.
\end{itemize}
\end{remark}

\begin{remark}
This separation leads to suboptimal solutions, as geometry and dynamics
are intrinsically coupled.
\end{remark}

\begin{remark}
The concept of \emph{Schwerpunkt-Trassierung} aims to unify these aspects
by defining a single design principle based on the motion of a
representative point.
\end{remark}

---

\section{Conceptual Idea}

\begin{definition}[Reference trajectory]
Let
\[
\gamma_c : I \to \R^2
\]
denote a reference trajectory representing the motion of a
characteristic point of the vehicle system.
\end{definition}

\begin{remark}
This point may be interpreted as:
\begin{itemize}
\item the center of mass,
\item a dynamically relevant reference point,
\item or an abstract design trajectory.
\end{itemize}
\end{remark}

---

\section{Relation to Track Geometry}

\begin{remark}
The physical track centerline \(\gamma(s)\) and cant angle \(\phi(s)\)
induce a spatial configuration of the vehicle.
\end{remark}

\begin{definition}[Induced reference trajectory]
Given a pose3 state
\[
(\gamma(s),\theta(s),\phi(s)),
\]
the induced reference trajectory is
\[
\gamma_c(s)
=
\gamma(s)
+
d\, n(s,\phi(s)),
\]
where:
\begin{itemize}
\item \(d\) is a characteristic offset,
\item \(n(s,\phi)\) is a direction depending on orientation and cant.
\end{itemize}
\end{definition}

\begin{remark}
The exact form of \(n(s,\phi)\) depends on the chosen mechanical model.
\end{remark}

\RESEARCH{Precise definition of offset direction based on vehicle model.}

---

\section{Schwerpunkt Principle}

\begin{definition}[Schwerpunkt principle]
A railway alignment is said to satisfy the Schwerpunkt principle if the
induced reference trajectory \(\gamma_c(s)\) possesses prescribed
smoothness and dynamical properties.
\end{definition}

\begin{remark}
In this formulation, the primary design object is not the track itself,
but the induced trajectory of the reference point.
\end{remark}

---

\section{Variational Formulation}

\begin{definition}[Schwerpunkt functional]
Let
\[
J_c[\gamma_c]
=
\int_0^L
\|\gamma_c''(s)\|^2 \,\dd s.
\]
\end{definition}

\begin{remark}
This functional penalizes curvature of the reference trajectory and
therefore promotes smooth motion.
\end{remark}

\begin{definition}[Schwerpunkt optimization problem]
Find
\[
(\kappa,\phi)
\]
such that the induced trajectory \(\gamma_c\) minimizes
\[
J_c[\gamma_c]
\]
subject to the pose3 dynamics.
\end{definition}

---

\section{Coupling to Pose3}

\begin{remark}
The trajectory \(\gamma_c\) depends on:
\[
\gamma(s), \quad \theta(s), \quad \phi(s).
\]
\end{remark}

\begin{remark}
Thus, the Schwerpunkt problem becomes a constrained optimization problem
in the variables:
\[
\kappa(s), \quad \phi(s).
\]
\end{remark}

---

\section{Interpretation}

\begin{remark}
The Schwerpunkt formulation shifts the perspective:
\begin{itemize}
\item from track geometry
\item to vehicle motion.
\end{itemize}
\end{remark}

\begin{remark}
This leads to a unified framework in which:
\begin{itemize}
\item geometry,
\item cant,
\item and dynamics
\end{itemize}
are optimized simultaneously.
\end{remark}

---

\section{Relation to Classical Design}

\begin{proposition}
Classical transition curves correspond to special cases of the
Schwerpunkt formulation where the reference trajectory coincides with
the track centerline.
\end{proposition}

\begin{remark}
In this case, the problem reduces to curvature-based optimization in
the plane.
\end{remark}

---

\section{Advantages}

\begin{itemize}
\item unified modelling of geometry and dynamics,
\item natural integration of cant,
\item compatibility with optimization frameworks,
\item extensibility to vehicle-specific models.
\end{itemize}

---

\section{Open Problems}

\OPEN{Precise mechanical interpretation of the reference trajectory.}

\OPEN{Integration with full vehicle dynamics models.}

\OPEN{Extension to network-level alignment design.}

---

\section{Connection to the AIM Framework}

\begin{summary}
Schwerpunkt-Trassierung provides a conceptual and mathematical extension
of classical alignment design.

Instead of prescribing geometric elements, it defines the design
problem in terms of a reference trajectory derived from the railway
state.

This aligns with the AIM philosophy of treating alignment as a
functional object governed by curvature and its derived quantities.
\end{summary}

\section{Vehicle Model}

\subsection{Rigid Body Approximation}

\begin{remark}
We model the railway vehicle as a rigid body with a characteristic
center of mass.
\end{remark}

\begin{definition}[Center of mass]
Let
\[
C(s) \in \R^3
\]
denote the spatial position of the vehicle center of mass.
\end{definition}

\begin{remark}
For planar track geometry, we consider the projection
\[
\gamma_c(s) \in \R^2
\]
of the center of mass trajectory.
\end{remark}

---

\subsection{Track Coordinate Frame}

\begin{definition}[Frenet frame]
Let
\[
t(s) = (\cos\theta(s),\sin\theta(s))
\]
be the tangent direction and
\[
n(s) = (-\sin\theta(s),\cos\theta(s))
\]
the normal direction.
\end{definition}

\begin{remark}
The pair \((t,n)\) defines a local coordinate system attached to the
track.
\end{remark}

---

\subsection{Cant Rotation}

\begin{definition}[Rotated normal]
Let the cant angle \(\phi(s)\) define a rotation of the cross-section.
\end{definition}

\begin{remark}
In a simplified model, the lateral offset direction remains aligned
with \(n(s)\), while vertical effects are captured implicitly.
\end{remark}

---

\subsection{Center of Mass Offset}

\begin{definition}[Offset model]
Let \(d>0\) denote a characteristic lateral offset.
The center of mass trajectory is approximated by
\[
\gamma_c(s)
=
\gamma(s) + d\,n(s).
\]
\end{definition}

\begin{remark}
More refined models may include dependence on \(\phi(s)\), e.g.
\[
\gamma_c(s)
=
\gamma(s) + d\,\cos\phi(s)\,n(s).
\]
\end{remark}

\RESEARCH{Inclusion of full 3D vehicle geometry.}

---

\section{Dynamics of the Center of Mass}

\subsection{Velocity and Acceleration}

\begin{proposition}
The velocity of the center of mass is
\[
\gamma_c'(s)
=
t(s) + d\,n'(s).
\]
\end{proposition}

\begin{remark}
Using
\[
n'(s) = -\kappa(s)\,t(s),
\]
we obtain
\[
\gamma_c'(s)
=
(1 - d\kappa(s))\,t(s).
\]
\end{remark}

---

\begin{proposition}
The acceleration of the center of mass is proportional to
\[
\gamma_c''(s)
=
- d\kappa'(s)\,t(s)
+
(1 - d\kappa(s))\kappa(s)\,n(s).
\]
\end{proposition}

\begin{remark}
Thus, both curvature and its derivative influence the motion of the
center of mass.
\end{remark}

---

\subsection{Effective Forces}

\begin{definition}[Lateral acceleration]
The lateral acceleration of the center of mass is
\[
a_c(s)
=
v^2 (1 - d\kappa(s))\kappa(s).
\]
\end{definition}

\begin{remark}
Including cant yields the effective acceleration
\[
a_{\mathrm{eff}}(s)
=
v^2 \kappa(s) - g\sin\phi(s).
\]
\end{remark}

---

\section{Coupled Design Problem}

\begin{definition}[Vehicle-based objective]
Define
\[
J_{\mathrm{veh}}[\kappa,\phi]
=
\int_0^L
\|\gamma_c''(s)\|^2
\,\dd s.
\]
\end{definition}

\begin{remark}
This functional directly measures the smoothness of the motion of the
center of mass.
\end{remark}

---

\begin{definition}[Dynamic objective]
Define
\[
J_{\mathrm{dyn}}[\kappa,\phi]
=
\int_0^L
\Bigl(
v^2 \kappa(s) - g\sin\phi(s)
\Bigr)^2
\,\dd s.
\]
\end{definition}

---

\begin{definition}[Combined vehicle functional]
A combined objective is
\[
J[\kappa,\phi]
=
\alpha J_{\mathrm{veh}}
+
\beta J_{\mathrm{dyn}}
+
\gamma \int_0^L (\kappa'(s))^2 \,\dd s
+
\delta \int_0^L (\phi'(s))^2 \,\dd s.
\]
\end{definition}

---

\section{Schwerpunkt-Trassierung (Formal Definition)}

\begin{definition}[Schwerpunkt alignment]
A Schwerpunkt alignment is a pair
\[
(\kappa,\phi)
\]
that minimizes the functional
\[
J[\kappa,\phi]
\]
subject to:
\begin{itemize}
\item pose3 dynamics,
\item boundary conditions,
\item engineering constraints.
\end{itemize}
\end{definition}

---

\section{Interpretation}

\begin{remark}
In this formulation:
\begin{itemize}
\item the track is no longer the primary design object,
\item the motion of the vehicle becomes central.
\end{itemize}
\end{remark}

\begin{remark}
Classical alignment design is recovered as a special case where
\[
d=0
\quad \text{and} \quad
\phi(s)=0.
\]
\end{remark}

---

\section{Discussion}

\begin{remark}
The presented model is intentionally simplified.
However, it already captures:
\begin{itemize}
\item coupling between curvature and cant,
\item influence of curvature derivatives,
\item relation between geometry and dynamics.
\end{itemize}
\end{remark}

\RESEARCH{Extension to full multibody vehicle models.}

\RESEARCH{Calibration using measured track data.}

\OPEN{Integration into practical design workflows.}

---

\section{Connection to the AIM Framework}

\begin{summary}
The combination of pose3 and a vehicle-based model provides a rigorous
foundation for Schwerpunkt-Trassierung.

It transforms alignment design into a problem of optimizing the motion
of a representative physical system.

This supports the central AIM idea that alignment is not merely a
geometric object, but a functional entity governed by curvature and
its physical implications.
\end{summary}
